\chapter*{Avant-propos}
\addcontentsline{toc}{chapter}{Avant-propos}

Ce document rassemble les corrigés de tous les exercices du cours
\og Méthodes numériques pour physiciens et ingénieurs \fg{} : les exercices
placés dans les chapitres I à VI, puis les dix-neuf exercices de la fin du
document.

\section{Comment chaque corrigé est construit}

Chaque corrigé suit le même plan, dans cet ordre :
\begin{enumerate}[label=\arabic*.]
  \item \textbf{Pourquoi ?} On explique d'abord l'idée : quel est le
        problème, pourquoi la méthode choisie convient, quelle difficulté
        elle résout et à quelles conditions elle marche.
  \item \textbf{Comment ? (étapes)} On déroule ensuite la solution étape par
        étape : mise en équations, calculs intermédiaires, puis algorithme et,
        quand l'énoncé le demande, programme.
  \item \textbf{Vérification numérique.} Chaque résultat chiffré et chaque
        algorithme ont été réellement calculés ou exécutés. Les valeurs
        données sont celles que ces calculs ont produites.
  \item \textbf{Références.} On indique où trouver la méthode dans le livre de
        référence ou dans une source publique.
\end{enumerate}
Lorsqu'un énoncé est incomplet, ambigu ou semble contenir une erreur, on le
dit dans un paragraphe \og Remarque sur l'énoncé \fg{}, et on précise
l'interprétation retenue. Rien n'est supposé sans être signalé.

\section{Références utilisées}

La référence principale est le livre
R. L. Burden, J. D. Faires et A. M. Burden, \textit{Numerical Analysis},
10\textsuperscript{e} édition, Cengage Learning, 2016, noté \textbf{[BF]}.
Les renvois du type \BF{alg.~6.1, p.~368} donnent le numéro de l'algorithme ou
du théorème et la page imprimée du livre. Les autres sources (articles,
encyclopédies en ligne) sont listées dans la bibliographie à la fin du
document.

\section{Programmes et calculs}

Les programmes demandés par les énoncés se trouvent dans le dossier
\fichier{latex-corriges/programmes/} (langages C, Fortran 90 et Pascal). Ils ont
été compilés avec \texttt{gcc}, \texttt{gfortran} et Free Pascal
(\texttt{fpc}) ; le script \fichier{tester.sh} les compile et les exécute
tous avec les données utilisées dans ce document (résultats dans
\fichier{calculs/programmes_resultats.txt}). Les calculs numériques et les
tests d'algorithmes sont dans \fichier{latex-corriges/calculs/} (Python avec
NumPy et SciPy) : \fichier{calculs.py}, \fichier{complements.py} et
\fichier{tests_algorithmes.py}, avec leurs sorties dans les fichiers
\fichier{*_resultats.txt} et \fichier{resultats.txt}.

Dans les algorithmes, la flèche $\leftarrow$ désigne l'affectation, et les
notations sont celles du cours.
